\section{INTRODUCTION}
\label{sec:introduction}

In eXtended Reality (XR), avatars serve as digital stand-ins for users, supporting natural interaction and enhancing the sense of presence.
By mirroring physical movements, avatars function not only as control interfaces but also as psychological proxies--entities that shape how users feel and behave within immersive environments.
The design of avatars--including their appearance, motion fidelity, and contextual responsiveness--affects user ownership, embodiment, and social presence~\cite{yun2023animationfidelity}. 
To be effective, avatars must replicate a user’s body accurately and responsively across diverse contexts.
This often requires adjustments to avatar shape and motion mapping--that is, the way sensor data drives the avatar rig--to accommodate interaction demands, device limitations, and environmental constraints~\cite{cheymol2023beyond}.
These requirements vary significantly depending on the available input devices.
For example, with only an HMD and hand controllers, the system must estimate the torso and legs movements, whereas with full-body trackers, it can rely on measured joint data and minimize estimation.

Recent advances in sparse-sensing have shown that full-body motion can be plausibly reconstructed from minimal inputs such as headsets, controllers, or commodity IMUs.
These methods are based on learning to infer missing kinematics from partial observations, often enhanced through spatial inference models such as transformers~\cite{jiang2022avatarposer,yang2024divatrack}.
Nonetheless, configuring and deploying personalized avatars across heterogeneous XR platforms remains technically complex and fragmented. 
Inconsistent coordinate systems and SDK-dependent retargeting pipelines, in particular, hinder seamless motion reproduction and cross-platform usability.
Although standard APIs and standards~\cite{khronos_openxr_1_0,openxr_ext_hand} facilitate access to tracking data, they do not provide a comprehensive end-to-end runtime environment.
Consequently, developers must still normalize heterogeneous inputs, route motion through configurable processing stages, and retarget outputs to diverse rigs under device variability.

To address these limitations, we propose \textbf{MotionDAG}, a modular motion integration framework that places motion processing at the core of XR runtime design.
MotionDAG addresses this gap with a modular runtime that links creator-facing node assets to an extensible developer codebase, so both audiences configure XR trackers and avatar rigs through the same directed acyclic graph (DAG). 
Rather than offering access to raw tracking data alone--as in standard APIs--MotionDAG provides a full runtime environment that standardizes inputs, composes motion processing, and aims to provide more predictable end-to-end behavior.

The system abstracts diverse coordinate spaces and tracking conventions into a standardized pose stream using a vendor-agnostic input layer.
An event-driven node-graph runtime composes motion operations into modular, reconfigurable pipelines, including inverse kinematics (IK) and regression-based estimation.
A key feature is the use of a typed \emph{Kinematic Buffer}, which enables efficient, zero-copy motion data sharing across nodes through a publish-subscribe model.
Data flows as a staged cascade rather than arbitrary node-to-node transfer, with nodes emitting typed outputs only to downstream stages.
Further, selected computationally intensive nodes can be wrapped and offloaded to remote servers without changing interface contracts in the rest of the graph.
At the output boundary, SDK-independent retargeting uses anatomical frame mappings. 
It supports expressive embodiment from partial hand models to full-body rigs.

By decoupling input normalization, motion estimation, and output retargeting into discrete, composable stages, MotionDAG supports flexible deployment models, in situ or partially ex situ, while maintaining interface consistency.
Together, these architectural principles support scalable development and deployment of custom avatars, reducing bespoke integration work while enabling consistent motion behavior under diverse XR configurations.

The remainder of this paper is organized as follows: \hyperref[sec:systemdesign]{System Design} describes the system architecture and the key design principles behind MotionDAG, including its node-graph processing model and abstraction strategies.
This is followed by a breakdown of key implementation elements in \hyperref[sec:implementation]{Runtime Implementation}.
In \hyperref[sec:workflow]{Configuring Workflow}, we detail the developer-facing workflow and show how the same configuration artifacts generalize across heterogeneous devices.
Further, \hyperref[sec:evaluation]{Evaluation} reports the controlled user study first, followed by the runtime and offloading analysis.
Finally, \hyperref[sec:discussionandconclusion]{Conclusion} discusses implications, limitations, and future directions.
